\clearpage
\begin{summarybox}
\begin{center}\begin{tikzpicture}[font=\small,>=stealth,line join=round]
\draw[gray!40,rounded corners=3pt] (0,0) rectangle (14.2,8.3);\node at (7.1,7.700000000000001) {顶点虽低，却不在允许的区间中};
\begin{scope}[shift={(1.8777777777777778,2.2249999999999996)},x=1.111111111111111cm,y=0.575cm]\draw[->,gray] (-0.7,0)--(4.7,0) node[right]{$x$};\draw[->,gray] (0,-1)--(0,7) node[above]{$y$};\draw[gray,thick,solid,domain=-0.2:4.2,samples=180] plot(\x,{(\x-2)^2+1});\draw[blue,thick,very thick,domain=3:4,samples=180] plot(\x,{(\x-2)^2+1});\fill[orange!85!black] (3,2) circle(2.2pt);\fill[orange!85!black] (4,5) circle(2.2pt);\draw[violet,dashed] (2,1.05)--(2,6.5);\draw[orange!85!black,very thick] (3,-.5)--(4,-.5);\node[below] at (2,1){$(2,1)$};\node[right] at (3,2){$(3,2)$};\node[above left=3pt] at (4,5){$(4,5)$};\foreach \v in {2,3,4}{\node[below=5pt] at (\v,-.5){$\v$};}\end{scope}\node[anchor=west,align=left,text width=5.0cm] at (8.4,4.0) {$f=(x-2)^2+1$\\[9pt]允许 $x\in[3,4]$\\[9pt]$2\notin[3,4]$\\[9pt]最近 $x=3$：最小值2\\[9pt]最远 $x=4$：最大值5};
\node[text width=12.6cm,align=center] at (7.1,.65) {灰线仅供定位；蓝色粗线是允许段，不能到灰色顶点处取值。};
\end{tikzpicture}\end{center}


\textbf{\textcolor{CSummary}{一句话抓住本质}}

\textbf{在一段允许的横坐标中，先找离对称轴最近和最远的地方；开口向上时近小远大，开口向下时近大远小。}

\medskip
\textbf{\textcolor{CSummary}{做题四步}}
\begin{enumerate}[leftmargin=*,itemsep=0.3em]
\item \textbf{配方或求轴：} 找到 $x=h=-\dfrac{b}{2a}$。
\item \textbf{看范围：} 标出闭区间 $[m,n]$，检查 $h$ 在左边、里面还是右边。
\item \textbf{看开口：} $a>0$ 近小远大；$a<0$ 近大远小。
\item \textbf{最后代入：} 最近处为 $h$（若在区间内）或最近端点；最远处在两端点中选。只算真正需要的值。
\end{enumerate}

\medskip
\end{summarybox}\clearpage\begin{summarybox}\textbf{\textcolor{CSummary}{加速小技巧}}

区间中点 $s=\dfrac{m+n}{2}$，用它来找\textbf{最远端点}：
\[
\boxed{h<s\Rightarrow n\text{ 更远};\quad h>s\Rightarrow m\text{ 更远};\quad h=s\Rightarrow\text{两端同高}.}
\]
注意“同高”说的是 $f(m)=f(n)$，并不是它们都等于顶点值。

\medskip
\textbf{\textcolor{CSummary}{一个最短的例子}}

$f(x)=(x-2)^2+1$，$x\in[3,4]$：轴 $x=2$ 在区间左边，最近处 $x=3$，最远处 $x=4$；又因为开口向上，故\textbf{最小值 $2$，最大值 $5$}。

\medskip
\textbf{\textcolor{CSummary}{两个必须检查的条件}}
\begin{itemize}[leftmargin=*,itemsep=0.3em]
\item 这是二次函数（$a\ne0$），且讨论的确实是一段\textbf{包含两端点的闭区间}。
\item \textbf{顶点只有在允许区间内才能取到}；离开了区间，就不能拿它的值充当区间最值。
\end{itemize}
\end{summarybox}
